\documentclass[11pt]{amsart}

\usepackage[T1]{fontenc}
\usepackage{lmodern}
\usepackage{mathtools,amssymb}
\usepackage{array,booktabs}
\usepackage[margin=0.85in]{geometry}
\usepackage{microtype}
\usepackage[hidelinks]{hyperref}

\newcommand{\Z}{\mathbb Z}
\newcommand{\B}{\mathcal B_{41}}
\newcommand{\A}{\mathcal A_{41}}
\newcommand{\PhiMap}{\Phi_{41}}

\setlength{\parindent}{0pt}
\setlength{\parskip}{0.45em}

\title[The Tunnell map at $n=41$]
{The Tunnell Map at $n=41$:\\
A First Walkthrough and the Complete Table}

\date{September 2026}

\hypersetup{
  pdftitle={The Tunnell Map at n=41: A First Walkthrough and the Complete Table},
  pdfauthor={}
}

\begin{document}

\maketitle

\section*{1. What is the problem?}

No number theory is needed to follow the calculations in this note.  We use
only integer triples, squares, parity, and divisibility.

An \emph{integer} is a whole number such as
\(-3,-2,-1,0,1,2,3\).  An \emph{integer triple} is an ordered list of three
integers, for example \((2,-1,2)\).  We begin with the equation
\[
 2x^2+y^2+8z^2=41. \tag{source equation}
\]
Every integer triple \((x,y,z)\) satisfying this equation is called a
\emph{source point}.  We write \(\B\) for the list of all source points.

We also consider the equation
\[
 2u^2+v^2+32w^2=41. \tag{target equation}
\]
Every integer triple \((u,v,w)\) satisfying this equation is called a
\emph{target point}.  We write \(\A\) for the list of all target points.
In standard set notation,
\[
\begin{aligned}
 \B&=\{(x,y,z)\in\Z^3:2x^2+y^2+8z^2=41\},\\
 \A&=\{(u,v,w)\in\Z^3:2u^2+v^2+32w^2=41\}.
\end{aligned}
\]
The symbol \(\Z^3\) simply means ``all ordered triples of integers.''

A direct finite check gives
\[
 |\B|=32,
 \qquad
 |\A|=16.
\]
Thus there are twice as many source points as target points.  The goal is not
merely to count them.  We want one fixed rule that sends every source point to
a target point, with exactly two source points landing on each target point.

\newpage
\section*{2. What does the map notation mean?}

The symbol
\[
 \PhiMap:\B\longrightarrow\A
\]
denotes our rule.  You may think of \(\PhiMap\) as a machine: a source point
is the input, and a target point is the output.  For example, one calculation
below will show that
\[
 \PhiMap(-4,-1,1)=(2,-1,1).
\]

If \(Y\) is a target point, a \emph{preimage} of \(Y\) is an input that the
machine sends to \(Y\).  The collection of all such inputs is called the
\emph{fibre} over \(Y\):
\[
 \PhiMap^{-1}(Y)=\{X\in\B:\PhiMap(X)=Y\}.
\]
The superscript \(-1\) here does not mean the number \(1/\PhiMap\).  It means
``look backward and collect all inputs that produce this output.''

The statement that \(\PhiMap\) is \emph{two-to-one} means
\[
 \text{every target point has exactly two preimages.}
\]
The two examples below calculate both preimages, not just one of them.

\begin{center}
\renewcommand{\arraystretch}{1.25}
\begin{tabular}{@{}>{\bfseries}p{0.17\textwidth}p{0.70\textwidth}@{}}
\toprule
Word & Meaning in this note \\
\midrule
source & An integer triple satisfying the source equation. \\
target & An integer triple satisfying the target equation. \\
map & A rule assigning one target to each source. \\
preimage & A source that maps to a specified target. \\
fibre & The complete collection of preimages of one target. \\
even/odd & An integer is even if it is divisible by \(2\); otherwise it is odd. \\
\(a\mid b\) & ``\(a\) divides \(b\)'': there is an integer \(k\) with \(b=ak\). \\
branch & One case of a rule defined by several different cases. \\
\bottomrule
\end{tabular}
\end{center}

\section*{3. The map in plain language}

The rule first looks at the third coordinate \(z\) of a source
\((x,y,z)\).

\subsection*{Case A: the third coordinate is even}

Write \(z=2w\), and halve the third coordinate:
\[
 (x,y,2w)\longmapsto(x,y,w). \tag{even rule}
\]
Why does this work?  If the source equation holds, then
\[
 2x^2+y^2+8(2w)^2
 =2x^2+y^2+32w^2=41,
\]
which is exactly the target equation.  Conversely, each target \((u,v,w)\)
has exactly one even source above it, namely \((u,v,2w)\).

This already gives one preimage of every target.  We still need one more.

\subsection*{Case B: the third coordinate is odd}

For an odd source \((x,y,z)\), we perform three divisibility tests:
\[
\begin{array}{c|c|c}
\text{test}&\text{condition}&\text{rule used if the condition holds}\\ \hline
1&8\mid(x+2z+y)&\Lambda_1\\
2&8\mid(y-x+2z)&\Lambda_2\\
3&4\mid x&\Lambda_3.
\end{array}
\]
The paper proves that at most one test can succeed for the same source
(Proposition~3.1).  If a test succeeds, the corresponding explicit formula
is used.  These are the \emph{direct} odd sources.  At \(n=41\), every direct
odd source uses the third formula
\[
 \Lambda_3(x,y,z)=\left(2z,y,-\frac{x}{4}\right). \tag{direct rule at \(n=41\)}
\]
The condition \(4\mid x\) ensures that \(-x/4\) is an integer.

If none of the three tests succeeds, the point is called \emph{residual}.
The residual points are paired with the remaining targets by the matching
procedure explained before Worked Example~2.

The two odd mechanisms are alternatives:
\[
 \boxed{\text{an odd source is either direct or residual, never both.}}
\]
The direct formulas pair their sources one-to-one with some targets.  The
residual matching pairs all remaining odd sources one-to-one with all
remaining targets.  Therefore every target receives one even preimage and
one odd preimage.  This is the structure behind the two-to-one result.

\paragraph{Where this appears in the paper.}
The even rule is equation~(4).  The three direct formulas are equations
(8)--(10), and their one-to-one property is Theorem~3.2.  The complete map is
Theorem~1.2.

\newpage
\section*{4. Worked Example 1: a direct fibre}

We will find \emph{every} source point sent to
\[
 Y=(2,-1,1).
\]

\subsection*{Step 1: Check the target point}

Insert \(u=2\), \(v=-1\), and \(w=1\) into the target equation:
\[
\begin{aligned}
 2u^2+v^2+32w^2
 &=2(2)^2+(-1)^2+32(1)^2\\
 &=2\cdot4+1+32\\
 &=8+1+32=41.
\end{aligned}
\]
Therefore \(Y\in\A\).

\subsection*{Step 2: Find the even preimage}

The even rule halves the source's third coordinate.  To run that rule
backward, we double the target's third coordinate:
\[
 (u,v,w)=(2,-1,1)
 \quad\leadsto\quad
 X_{\mathrm{even}}=(u,v,2w)=(2,-1,2).
\]
Check that this is a source:
\[
 2(2)^2+(-1)^2+8(2)^2=8+1+32=41.
\]
Applying the map now gives
\[
 \PhiMap(2,-1,2)=(2,-1,1)=Y.
\]

\subsection*{Step 3: Find the odd preimage}

For the third direct branch, the inverse formula in Theorem~3.2 is
\[
 (u,v,w)\longmapsto\left(-4w,v,\frac{u}{2}\right)
 \quad\text{when }u\equiv2\pmod4.
\]
The notation \(u\equiv2\pmod4\) means that \(u-2\) is divisible by \(4\).
Here \(u=2\), so the condition holds.  Substitution gives
\[
 X_{\mathrm{odd}}
 =\left(-4(1),-1,\frac{2}{2}\right)
 =(-4,-1,1).
\]
Its third coordinate is odd.  Check that it is a source:
\[
 2(-4)^2+(-1)^2+8(1)^2=32+1+8=41.
\]

\newpage
\subsection*{Step 4: Check which odd rule applies}

For \((x,y,z)=(-4,-1,1)\), the three tests become
\[
\begin{array}{c|c|c}
\text{test}&\text{value}&\text{does the required number divide it?}\\ \hline
1&x+2z+y=-4+2-1=-3&8\nmid-3\\
2&y-x+2z=-1+4+2=5&8\nmid5\\
3&x=-4&4\mid-4.
\end{array}
\]
Thus only Test~3 succeeds, so \(\Lambda_3\) is the correct branch.

\subsection*{Step 5: Apply the formula}

Substitute \(x=-4\), \(y=-1\), and \(z=1\):
\[
\begin{aligned}
 \PhiMap(-4,-1,1)
 &=\Lambda_3(-4,-1,1)\\
 &=\left(2(1),-1,-\frac{-4}{4}\right)\\
 &=(2,-1,1)=Y.
\end{aligned}
\]

\subsection*{Step 6: Collect the complete answer}

We found one even preimage and one odd preimage:
\[
 \boxed{\PhiMap^{-1}(2,-1,1)
 =\{(2,-1,2),\,(-4,-1,1)\}.}
\]
There cannot be another even preimage because doubling \(w\) gives only one
integer.  There cannot be another odd preimage because Theorem~3.2 proves
that the direct rule is one-to-one onto its target set; residual sources are
assigned only to targets not already used by the direct rule.  Theorem~1.2
assembles these facts and states that the displayed list is the whole fibre.

\paragraph{Paper references.}
Even rule: equation~(4).  Branch test: Proposition~3.1.  Formula:
equation~(10).  Direct bijection: Theorem~3.2.  Complete fibre:
Theorem~1.2.

\newpage
\section*{5. The extra idea needed for a residual source}

Worked Example~1 used a short formula.  Some odd source points fail all three
direct tests.  The second example shows what happens then.

\subsection*{Why points are first grouped into sign-pairs}

If \((a,b,c)\) satisfies either of our equations, then
\((-a,-b,-c)\) also satisfies it, because squaring removes every minus sign.
The residual construction temporarily treats these two points as one
\emph{sign-pair}:
\[
 \{(a,b,c),(-a,-b,-c)\}.
\]
The paper calls such a sign-pair an \emph{antipodal orbit}.  In this note,
``orbit'' means only this two-element pair.

\subsection*{Why the third coordinate is rescaled}

For residual bookkeeping, the paper replaces
\[
 (x,y,z)\longmapsto(x,y,2z)
 \quad\text{and}\quad
 (u,v,w)\longmapsto(u,v,4w).
\]
These are called \emph{lifted coordinates}.  They put both sides into the
same equation
\[
 Q(a,b,c)=2a^2+b^2+2c^2=41.
\]
Nothing mysterious happens here: to return to the original coordinates, we
halve the source's last coordinate or divide the target's last coordinate by
four.

After the direct points have been removed, the four residual source
sign-pairs and four residual target sign-pairs at \(n=41\) are represented by
\[
\begin{array}{c|c@{\qquad}c|c}
\text{source label}&\text{listed source}&\text{target label}&\text{listed target}\\
\midrule
s_1&(-2,-5,-2)&t_1&(-4,-3,0)\\
s_2&(-2,-5, 2)&t_2&(-4, 3,0)\\
s_3&(-2, 5,-2)&t_3&(0,-3,-4)\\
s_4&(-2, 5, 2)&t_4&(0,-3, 4).
\end{array}
\]
Each listed point stands for itself together with its negative.

\subsection*{Where one of the keys comes from}

The keys are arithmetic data, not arbitrary preference numbers.  Here is one
complete calculation.  Take the listed points
\[
 E=(-2,-5,2)\quad\text{from }s_2,
 \qquad
 F=(-4,-3,0)\quad\text{from }t_1.
\]
Form the two midpoint vectors
\[
 \frac{E+F}{2}=(-3,-4,1),
 \qquad
 \frac{E-F}{2}=(1,-1,1).
\]
The paper changes a vector's sign, if needed, so that its first nonzero
coordinate is positive.  Thus the first direction becomes \((3,4,-1)\),
while the second stays \((1,-1,1)\).  If all three coordinates had a common
factor, we would first divide by their greatest common divisor.  No division
is needed here.

For a direction \(d=(d_1,d_2,d_3)\), its key is
\[
 \bigl(Q(d),d_1,d_2,d_3\bigr),
 \qquad
 Q(d)=2d_1^2+d_2^2+2d_3^2.
\]
The two keys in this example begin with
\[
 Q(3,4,-1)=36,
 \qquad
 Q(1,-1,1)=5.
\]
The smaller key is therefore
\[
 (5,1,-1,1).
\]
This is exactly the key assigned to the pair \(s_2t_1\) below.  The other
keys are obtained by the same calculation.

\subsection*{The matching procedure}

The construction assigns a four-number \emph{key} to every possible
source--target sign-pair.  Keys are compared from left to right; at the first
place where they differ, the smaller number wins.  For example,
\[
 (5,1,-1,-1)<(5,1,-1,1)
\]
because the first three entries agree and \(-1<1\) in the fourth entry.

Each source proposes to targets in increasing key order.  A target with no
current proposal holds the proposal.  If a target already holds one, it keeps
the proposal with the smaller key and rejects the other source.  A rejected
source tries its next target.  When every source is held by a different
target, the process stops.  The word \emph{stable} means that there is no
unmatched source--target pair that would both prefer one another to their
assigned partners.

For \(s_2\), the target order is
\[
 t_1,\quad t_4,\quad t_2,\quad t_3,
\]
with keys
\[
 (5,1,-1,1),\quad (5,1,1,1),\quad
 (20,1,-4,1),\quad (20,1,4,1).
\]
This is the row of data used in Worked Example~2.

\section*{6. Worked Example 2: a residual fibre}

We will find every source point sent to
\[
 Y=(4,-3,0).
\]

\subsection*{Step 1: Check the target point}

Substitute \(u=4\), \(v=-3\), and \(w=0\):
\[
 2(4)^2+(-3)^2+32(0)^2=32+9+0=41.
\]
Therefore \(Y\in\A\).

\subsection*{Step 2: Find the even preimage}

Double the last coordinate of the target:
\[
 X_{\mathrm{even}}=(u,v,2w)=(4,-3,0).
\]
Because \(w=0\), the triple happens to look unchanged.  It satisfies the
source equation:
\[
 2(4)^2+(-3)^2+8(0)^2=32+9=41.
\]
The even rule gives
\[
 \PhiMap(4,-3,0)=(4,-3,0)=Y.
\]

\subsection*{Step 3: Locate the target sign-pair}

The target lift multiplies the last coordinate by four.  Since our target has
last coordinate zero, its lift is still
\[
 (4,-3,0).
\]
The listed representative of \(t_2\) is \(F=(-4,3,0)\).  Its negative is
\[
 -F=(4,-3,0)=Y.
\]
Therefore our target belongs to the sign-pair \(t_2\).  To find its odd
preimage, we now determine which source sign-pair the matching assigns to
\(t_2\).

\subsection*{Step 4: Run the small matching for n=41}

The complete proposal run is shown below.  ``Held'' means tentatively
accepted; a held proposal can still be replaced by a better one.

\begin{center}
\small
\renewcommand{\arraystretch}{1.24}
\begin{tabular}{@{}c p{0.63\textwidth}@{}}
\toprule
proposal & what happens \\
\midrule
\(s_1\to t_1\) & \(t_1\) is empty, so it holds \(s_1\). \\
\(s_2\to t_1\) & Rejected.  The key for \(s_1t_1\) is
\((5,1,-1,-1)\), smaller than the key \((5,1,-1,1)\) for \(s_2t_1\). \\
\(s_3\to t_4\) & \(t_4\) is empty, so it holds \(s_3\). \\
\(s_4\to t_3\) & \(t_3\) is empty, so it holds \(s_4\). \\
\(s_2\to t_4\) & Rejected.  The key for \(s_3t_4\) is
\((5,1,-1,-1)\), smaller than the key \((5,1,1,1)\) for \(s_2t_4\). \\
\(s_2\to t_2\) & \(t_2\) is empty, so it holds \(s_2\).  All four sources
are now held, so the algorithm stops. \\
\bottomrule
\end{tabular}
\end{center}

The final pairing of sign-pairs is
\[
 s_1\leftrightarrow t_1,\qquad
 s_2\leftrightarrow t_2,\qquad
 s_3\leftrightarrow t_4,\qquad
 s_4\leftrightarrow t_3.
\]
In particular, \(t_2\) is matched with \(s_2\).  The listed representative of
\(s_2\) is
\[
 E=(-2,-5,2).
\]
Undoing the source lift means halving the last coordinate, which gives the
odd source candidate
\[
 X_{\mathrm{odd}}=(-2,-5,1).
\]

\newpage
\subsection*{Step 5: Verify the source and decide its target sign}

First check that the candidate is a source:
\[
 2(-2)^2+(-5)^2+8(1)^2=8+25+8=41.
\]
Its three direct tests are
\[
\begin{array}{c|c|c}
\text{test}&\text{value}&\text{result}\\ \hline
1&x+2z+y=-2+2-5=-5&8\nmid-5\\
2&y-x+2z=-5+2+2=-1&8\nmid-1\\
3&x=-2&4\nmid-2.
\end{array}
\]
All three tests fail, confirming that this source belongs to the residual
rather than the direct branch.

Matching the two sign-pairs is not yet enough: the actual source must map to
either \(F\) or its negative \(-F=(4,-3,0)\).  The construction compares the
two midpoint directions.

For \(E=(-2,-5,2)\) and \(F=(-4,3,0)\), the two half-sums are
\[
 \frac{E+F}{2}=(-3,-1,1),
 \qquad
 \frac{E-F}{2}=(1,-4,1).
\]
The first vector is the midpoint of \(E\) and \(F\).  The second is the
midpoint of \(E\) and \(-F\), because
\[
 \frac{E+(-F)}{2}=\frac{E-F}{2}.
\]
Changing the sign of the first vector gives the oriented direction
\((3,1,-1)\).  Their \(Q\)-values are
\[
 Q(3,1,-1)=2(3)^2+1^2+2(-1)^2=21,
\]
and
\[
 Q(1,-4,1)=2(1)^2+(-4)^2+2(1)^2=20.
\]
Since \(20<21\), the second choice wins.  It corresponds to \(-F\), so
\[
 E\longmapsto-F=(4,-3,0).
\]
Returning from the lifted source \(E=(-2,-5,2)\) to the original source
halves the last coordinate and gives \((-2,-5,1)\).  Hence
\[
 \PhiMap(-2,-5,1)=(4,-3,0)=Y.
\]

\subsection*{Step 6: Collect the complete answer}

Together with the even preimage from Step~2, we have
\[
 \boxed{\PhiMap^{-1}(4,-3,0)
 =\{(4,-3,0),\,(-2,-5,1)\}.}
\]
The residual algorithm pairs each remaining source sign-pair with exactly one
remaining target sign-pair.  Proposition~4.5 proves that this matching is the
unique stable one, equation~(12) supplies the point-level sign, and
Theorem~7.3 proves that Algorithm~7.1 computes the matching.  Theorem~1.2 then
shows that no third preimage exists.

\paragraph{Paper references.}
Sign lift: equation~(12).  Unique matching: Proposition~4.5.  Algorithm:
Algorithm~7.1.  Correctness of the algorithm: Theorem~7.3.  This complete
\(n=41\) run: Example~8.1.  Complete fibre: Theorem~1.2.

\newpage
\section*{7. The complete answer for n=41}

The next two tables contain every point.  Read one row horizontally:
\[
 \boxed{\text{target}\quad\longleftarrow\quad
 \text{its even preimage and its odd preimage}.}
\]
The first table contains the eight targets whose odd preimage uses the direct
formula.  The second contains the eight targets whose odd preimage comes from
the residual matching.

\subsection*{Direct fibres}

\begin{center}
\small
\renewcommand{\arraystretch}{1.22}
\begin{tabular}{@{}c c c c@{}}
\toprule
target & even preimage & odd preimage & odd rule \\
\midrule
$(-2,-1,-1)$ & $(-2,-1,-2)$ & $(4,-1,-1)$
  & $\Lambda_3(4,-1,-1)$ \\
$(-2,-1, 1)$ & $(-2,-1, 2)$ & $(-4,-1,-1)$
  & $\Lambda_3(-4,-1,-1)$ \\
$(-2, 1,-1)$ & $(-2, 1,-2)$ & $(4, 1,-1)$
  & $\Lambda_3(4,1,-1)$ \\
$(-2, 1, 1)$ & $(-2, 1, 2)$ & $(-4,1,-1)$
  & $\Lambda_3(-4,1,-1)$ \\
$( 2,-1,-1)$ & $( 2,-1,-2)$ & $(4,-1,1)$
  & $\Lambda_3(4,-1,1)$ \\
$( 2,-1, 1)$ & $( 2,-1, 2)$ & $(-4,-1,1)$
  & $\Lambda_3(-4,-1,1)$ \\
$( 2, 1,-1)$ & $( 2, 1,-2)$ & $(4,1,1)$
  & $\Lambda_3(4,1,1)$ \\
$( 2, 1, 1)$ & $( 2, 1, 2)$ & $(-4,1,1)$
  & $\Lambda_3(-4,1,1)$ \\
\bottomrule
\end{tabular}
\end{center}

For each row, applying \(\Lambda_3\) to the odd preimage gives the target in
the first column.  The sixth row is Worked Example~1.

\subsection*{Residual fibres}

\begin{center}
\small
\renewcommand{\arraystretch}{1.22}
\begin{tabular}{@{}c c c c@{}}
\toprule
target & even preimage & odd preimage & matching origin \\
\midrule
$(-4,-3,0)$ & $(-4,-3,0)$ & $(2,5,1)$
  & negative of the $s_1$ pair \\
$(-4, 3,0)$ & $(-4, 3,0)$ & $(2,5,-1)$
  & negative of the $s_2$ pair \\
$(0,-3,-1)$ & $(0,-3,-2)$ & $(-2,5,1)$
  & $s_4$ pair \\
$(0,-3, 1)$ & $(0,-3, 2)$ & $(-2,5,-1)$
  & $s_3$ pair \\
$(0, 3,-1)$ & $(0, 3,-2)$ & $(2,-5,1)$
  & negative of the $s_3$ pair \\
$(0, 3, 1)$ & $(0, 3, 2)$ & $(2,-5,-1)$
  & negative of the $s_4$ pair \\
$(4,-3,0)$ & $(4,-3,0)$ & $(-2,-5,1)$
  & $s_2$ pair \\
$(4, 3,0)$ & $(4, 3,0)$ & $(-2,-5,-1)$
  & $s_1$ pair \\
\bottomrule
\end{tabular}
\end{center}

The seventh row is Worked Example~2.

\newpage
\section*{8. Why this table is complete}

The direct table has eight different targets, and the residual table has
eight more.  Thus all \(16\) target points occur exactly once as table rows.
Each row contains two source points, one even and one odd, giving
\(16\cdot2=32\) source entries.  These entries are all different, so every
one of the \(32\) source points occurs exactly once.

Consequently, the tables show concretely that every target has exactly two
preimages.  They are not evidence from a few selected examples: together the
two tables are the full finite calculation at \(n=41\).

As a separate machine check, the table was regenerated from the executable
Lean definition
\[
 \texttt{tunnellMapTotal41}.
\]
The Lean theorems
\texttt{card\_BRep\_41} and \texttt{card\_ARep\_41} verify the counts
\(32\) and \(16\).  The theorem named
\[
 \texttt{tunnellMapTotal41\_exactly\_two}
\]
verifies that every target has two distinct preimages and that there are no
others.

\end{document}
